\documentclass[crop,tikz,11pt]{standalone}
\usepackage{geometricfigures}
\usepackage{amsmath,amssymb}
\usetikzlibrary{arrows.meta,patterns,calc}

\definecolor{rot}{HTML}{C1701E}
\definecolor{net}{HTML}{389826}    % the cap, filled net!35

% What fills the hole of the latent annulus that the region D = {H < H0, p < 0} (hatched) becomes,
% between the rotating orbit of energy H0 = 1.6 with p < 0 (orange) and the zero section nu (dotted);
% from the paper on symplectic autoencoders, arXiv:2312.10004.
%   (a), (b) the round chart r^2/2 = a - p (a = 3.2): D and the band 0 < p < a (grey), bent onto
%       the paraboloid of cylinder-to-poloidal; the band is closed by the apex (open circle), which
%       is not on M, so the origin of the latent plane is missing. Latent areas J and 2 pi a:
%       |int gamma| = J + 2 pi a, while the orbit has the action J.
%   (c), (d) a symplectic decoder of all of R^2: D closed by a cap C (green) attached along nu, on
%       the cylinder of cylinder-to-poloidal (a); the rest of M is left out. The latent curve gamma
%       encloses the area J = |preimage of D| (the action of the orbit), so the latent disk mapped
%       onto C has the area int_C omega = oint_nu p dtheta = 0. Panel (d) is drawn with r^2/2 = -p
%       for D, so that gamma encloses the area J; the green disk inside it is what the decoder would
%       need and has no room for.
% (b) and (d) are drawn to the same scale. Geometry as in cylinder-to-poloidal: a point (theta, p)
% at (x', y', p), x' = -r cos(theta+beta), y' = -r sin(theta+beta), orthographic from elevation e:
% page (X, Y) = (x', hc p + se y'); theta in (30, 210) faces the viewer.
\pgfmathsetmacro{\a}{3.2}
\pgfmathsetmacro{\Ho}{1.6}
\pgfmathsetmacro{\pc}{2.6}
\pgfmathsetmacro{\bet}{-30}
\pgfmathsetmacro{\se}{0.42}
\pgfmathsetmacro{\hc}{0.85}
\pgfmathsetmacro{\Rc}{sqrt(2*\a)}       % cylinder radius, = radius of phi(nu)
\pgfmathsetmacro{\Rb}{sqrt(2*(\a+\pc))} % bottom rim of the dome
\pgfmathsetmacro{\thr}{\se/\hc}        % visibility threshold on the dome
\pgfmathsetmacro{\pap}{\a-\thr*\thr/2} % top of the dome's silhouette
\pgfmathsetmacro{\Hcap}{2.2}            % height of the cap on the page
% p on the rotating orbit of energy H0 with p < 0, the outer boundary of D
\pgfmathdeclarefunction{po}{1}{\pgfmathparse{-sqrt(2*(\Ho-cos(#1)))}}
% a point (theta, p) with the circle at height p of radius r(p), r^2/2 = #2 - p (latent planes,
% dome) or of radius Rc (#2 < 0: the cylinder)
\pgfmathdeclarefunction{rof}{2}{\pgfmathparse{(#2) < 0 ? \Rc : sqrt(2*max(0,(#2)-(#1)))}}
\pgfmathdeclarefunction{lx}{3}{\pgfmathparse{-rof(#3,#2)*cos((#1)+\bet)}}
\pgfmathdeclarefunction{ly}{3}{\pgfmathparse{-rof(#3,#2)*sin((#1)+\bet)}}
\pgfmathdeclarefunction{pY}{3}{\pgfmathparse{\hc*(#3)+\se*ly(#1,#2,#3)}}

\begin{document}
\begin{tikzpicture}[>=Stealth,x=0.475cm,y=0.475cm,
  hidden/.style={opacity=0.38},
  hatchD/.style={pattern=north east lines,pattern color=fg!60!bg},
  lab/.style={font=\footnotesize,inner sep=1pt},
  mask/.style={fill=bg,fill opacity=0.85,text opacity=1}]

% \curves{A}: the orbit and nu on a 3d panel, and \hatchD{A}: D between them,
% front (theta in 30..210) solid, back faded; A = a on the dome, -1 on the cylinder
\newcommand{\curves}[1]{%
  \draw[rot,thick,hidden] plot[domain=210:390,samples=90,smooth] ({lx(\x,#1,po(\x))},{pY(\x,#1,po(\x))});
  \draw[fg!70!bg,densely dotted,thick,hidden] plot[domain=210:390,samples=90,smooth] ({lx(\x,#1,0)},{pY(\x,#1,0)});
  \draw[rot,thick] plot[domain=30:210,samples=90,smooth] ({lx(\x,#1,po(\x))},{pY(\x,#1,po(\x))});
  \draw[fg!70!bg,densely dotted,thick] plot[domain=30:210,samples=90,smooth] ({lx(\x,#1,0)},{pY(\x,#1,0)});}
\newcommand{\hatchD}[1]{%
  \fill[hatchD]
    plot[domain=30:210,samples=90,smooth] ({lx(\x,#1,po(\x))},{pY(\x,#1,po(\x))})
    -- plot[domain=210:30,samples=90,smooth] ({lx(\x,#1,0)},{pY(\x,#1,0)}) -- cycle;}

% ================================================================ (a) case (iii): the paraboloid
\fill[fg!4!bg]
  plot[domain=-\pc:\pap,samples=60,smooth] ({sqrt(max(0,2*(\a-\x)-\thr*\thr))},{\hc*\x+\se*\thr})
  -- plot[domain=\pap:-\pc,samples=60,smooth] ({-sqrt(max(0,2*(\a-\x)-\thr*\thr))},{\hc*\x+\se*\thr})
  -- plot[domain={180+asin(\thr/\Rb)-\bet}:{-asin(\thr/\Rb)-\bet},samples=60,smooth]
       ({lx(\x,\a,-\pc)},{pY(\x,\a,-\pc)}) -- cycle;
% the band 0 < p < a: above the front of nu, up to the silhouette
\fill[fg!16!bg]
  plot[domain={180+asin(\thr/\Rc)-\bet}:{-asin(\thr/\Rc)-\bet},samples=90,smooth] ({lx(\x,\a,0)},{pY(\x,\a,0)})
  -- plot[domain=0:\pap,samples=40,smooth] ({sqrt(max(0,2*(\a-\x)-\thr*\thr))},{\hc*\x+\se*\thr})
  -- plot[domain=\pap:0,samples=40,smooth] ({-sqrt(max(0,2*(\a-\x)-\thr*\thr))},{\hc*\x+\se*\thr}) -- cycle;
\hatchD{\a}
\draw[fg!25!bg]
  plot[domain=-\pc:\pap,samples=60,smooth] ({sqrt(max(0,2*(\a-\x)-\thr*\thr))},{\hc*\x+\se*\thr})
  -- plot[domain=\pap:-\pc,samples=60,smooth] ({-sqrt(max(0,2*(\a-\x)-\thr*\thr))},{\hc*\x+\se*\thr});
\draw[fg!30!bg,densely dotted] plot[domain=0:360,samples=120,smooth] ({lx(\x,\a,-\pc)},{pY(\x,\a,-\pc)});
\curves{\a}
\draw[fg,fill=bg,line width=0.5pt] (0,{\hc*\a}) circle (2.2pt);
\draw[fg!50!bg] (0.1,{\hc*\a+0.1}) -- (1.6,{\hc*\a+1.1})
  node[lab,fg,anchor=south west,align=left] {apex: the circle\\$p_\theta=a$};
\node[lab,mask] at (0,{\hc*1.4}) {$2\pi a$};
\node[lab,mask] at ({lx(140,\a,0.5*po(140))},{pY(140,\a,0.5*po(140))}) {$D$: $J$};
\node[lab,anchor=east,mask] at ({lx(45,\a,0)-0.3},{pY(40,\a,0)-0.1}) {$\nu$};

% ================================================================ (b) case (iii): the punctured plane
\begin{scope}[shift={(12.4,0)}]
  \fill[hatchD,even odd rule]
    plot[domain=0:360,samples=181,smooth] ({lx(\x,\a,po(\x))},{ly(\x,\a,po(\x))}) -- cycle
    (0,0) circle ({\Rc});
  \fill[fg!16!bg] (0,0) circle ({\Rc});
  \draw[fg!70!bg,densely dotted,thick] (0,0) circle ({\Rc});
  \draw[rot,thick] plot[domain=0:360,samples=181,smooth] ({lx(\x,\a,po(\x))},{ly(\x,\a,po(\x))}) -- cycle;
  \draw[fg,fill=bg,line width=0.5pt] (0,0) circle (2.2pt);
  \node[lab,anchor=north] at (0,-0.2) {$2\pi a$};
  \node[lab,anchor=south] at (0,{\Rc-1.0}) {$\varphi(\nu)$};
  \node[lab,rot!70!fg,anchor=south east] at ({lx(345,\a,po(345))},{ly(345,\a,po(345))}) {$\gamma$};
  \draw[fg!50!bg] ({lx(150,\a,0.5*po(150))},{ly(150,\a,0.5*po(150))}) -- (3.6,-2.9)
    node[lab,fg,anchor=north west] {$\varphi(D)$: $J$};
  \node[lab,anchor=north] at (0,-3.7) {$\lvert\operatorname{int}\gamma\rvert=J+2\pi a$};
\end{scope}
\draw[->,thick,fg!60!bg] (4.9,0) -- (8.6,0) node[midway,above,lab,fg] {$\varphi$};

% ================================================================ (c) case (iv): D and a cap
\begin{scope}[shift={(0,-10.5)}]
  \fill[fg!4!bg]
    plot[domain=30:210,samples=60,smooth] ({lx(\x,-1,-\pc)},{pY(\x,-1,-\pc)})
    -- plot[domain=210:390,samples=60,smooth] ({lx(\x,-1,0)},{pY(\x,-1,0)}) -- cycle;
  % the cap over nu: an elliptic dome from the back of nu, closed by the front of nu
  \fill[net!35!bg] ({-\Rc},0) arc[start angle=180,end angle=0,x radius=\Rc,y radius=\Hcap]
    -- plot[domain=210:30,samples=90,smooth] ({lx(\x,-1,0)},{pY(\x,-1,0)}) -- cycle;
  \draw[net!80!fg] ({-\Rc},0) arc[start angle=180,end angle=0,x radius=\Rc,y radius=\Hcap];
  \hatchD{-1}
  \draw[fg!25!bg] ({-\Rc},{-\hc*\pc}) -- ({-\Rc},0);
  \draw[fg!25!bg] ({\Rc},{-\hc*\pc}) -- ({\Rc},0);
  \draw[fg!30!bg,densely dotted] plot[domain=0:360,samples=120,smooth] ({lx(\x,-1,-\pc)},{pY(\x,-1,-\pc)});
  \curves{-1}
  \node[lab,anchor=south] at (0,{\Hcap+0.1}) {$C$: $\oint_\nu p_\theta\,d\theta=0$};
  \node[lab,mask] at ({lx(140,-1,0.5*po(140))},{pY(140,-1,0.5*po(140))}) {$D$: $J$};
  \node[lab,anchor=east,mask] at ({lx(45,-1,0)-0.3},{pY(40,-1,0)-0.1}) {$\nu$};
\end{scope}

% ================================================================ (d) case (iv): the latent plane R^2
\begin{scope}[shift={(12.4,-10.5)}]
  \fill[hatchD] plot[domain=0:360,samples=181,smooth] ({lx(\x,0,po(\x))},{ly(\x,0,po(\x))}) -- cycle;
  \fill[net!35!bg] (0,0) circle (0.75);
  \draw[net!80!fg] (0,0) circle (0.75);
  \draw[rot,thick] plot[domain=0:360,samples=181,smooth] ({lx(\x,0,po(\x))},{ly(\x,0,po(\x))}) -- cycle;
  \node[lab,rot!70!fg,anchor=south east] at ({lx(345,0,po(345))},{ly(345,0,po(345))}) {$\gamma$};
  \draw[fg!50!bg] (0.4,-0.4) -- (3.0,-1.6)
    node[lab,fg,anchor=west,align=left] {$(\Psi^\mathrm{dec})^{-1}(C)$:\\$J-J=0$};
  \draw[fg!50!bg] ({lx(150,0,0.6*po(150))},{ly(150,0,0.6*po(150))}) -- (3.0,0.9)
    node[lab,fg,anchor=west] {$(\Psi^\mathrm{dec})^{-1}(D)$: $J$};
  \node[lab,anchor=north] at (0,-3.2) {$\lvert\operatorname{int}\gamma\rvert=J$};
\end{scope}
\draw[<-,thick,fg!60!bg] (4.9,-10.5) -- (8.6,-10.5) node[midway,above,lab,fg] {$\Psi^\mathrm{dec}$};

\node[font=\footnotesize,anchor=north] at (0,-4.6) {\textbf{(a)}};
\node[font=\footnotesize,anchor=north] at (12.4,-4.6) {\textbf{(b)}};
\node[font=\footnotesize,anchor=north] at (0,-14.3) {\textbf{(c)}};
\node[font=\footnotesize,anchor=north] at (12.4,-14.3) {\textbf{(d)}};
\end{tikzpicture}
\end{document}
